\documentclass[12pt,a4paper]{article}
\usepackage{amsmath, amssymb, amsthm}
\usepackage{geometry}
\geometry{margin=1in}
\usepackage{hyperref}
\usepackage{enumitem}
\setlist{nosep}

\title{Fixed-Income Securities: Concise Summary}
\author{Navid}
\date{\today}

\begin{document}

\setlength{\parindent}{0pt}

\maketitle
\tableofcontents


%-------------------


\section{Credit Risk Fundamentals}

\subsection{Key Concepts and Stylized Facts}

\subsubsection{Credit Risk and Credit Spreads}

\textbf{Credit risk} is the possibility that a counterparty fails to meet contractual obligations, causing \textbf{price fluctuations that reflect this exposure to credit risk}. Exposure has two components: \textbf{timing of default} (measured by \textbf{default probability}) and \textbf{magnitude of default} (measured by \textbf{loss upon default}). Loss upon default can be decomposed into the product of the credit exposure (i.e., 
the transaction amount that is at risk), and the proportional loss given default.

Credit risk exposure on fixed-income securities is compensated by an additional yield called the \textbf{credit spread}. For a debt claim with value $D$ promising $M$ at period end: 

$$D = \frac{M}{1 + y} \quad \text{(8.1)}$$

If defaultable with EMM probability of default $q$ which results in delivering $M(1-L)$, where $L$ is \textbf{loss given default}:

$$D = \frac{M(1-q) + M(1-L)q}{1 + r} \quad \text{(8.2)}$$

where $r$ is the risk-free rate. Combining (8.1) and (8.2):

$$y - r \approx qL \quad \text{(8.3)}$$

\textbf{Fundamental decomposition:} credit spread = (default probability under EMM) $\times$ (loss given default).

\textbf{Total Yield spread vs. credit spread:} Yield spread of defaultable instruments includes credit spread plus other premiums (tax, liquidity); the fraction of the yield spread that can be attributed to credit risk is an empirical issue.

\subsubsection{The Term Structure of Credit Spreads}

Consider homogeneous defaultable securities (comparable credit risk source) with observable market prices at different maturities. Assuming credit risk is the only factor in total yield spread, the \textbf{term structure of credit spreads (TSCS)} is calculated as the difference between yield-to-maturity and (risk-free) zero-coupon yield for the same horizon.

\textbf{TSCS dynamics:} Credit spreads are volatile and counter-cyclical (spikes during 2007-08 crisis and 2020 Covid pandemic). Like risk-free zero-coupon yields, the TSCS is usually increasing during quiet times and inverted in crisis times.

The three typical shapes are increasing, decreasing or humped. It is therefore possible to smooth 
the TSCS using the Nelson-Siegel-Svensson.

Building a TSCS requires a homogeneous set of defaultable claims sharing similar credit risk exposure (comparable default likelihood and losses given default). This homogeneity is achieved by grouping claims by type of instrument (e.g., senior unsecured bonds), risk class (e.g., rating), or industry. However, stricter grouping criteria reduce the population size, potentially introducing noisier data.

\subsubsection{Recovery Rates}

\textbf{Recovery rate:} One minus loss given default; one of two major constituents of credit spread. For defaultable zero-coupon bond value $D(t,T)$ (promising \$1 at $T$), let $\tau$ = random default time and $\delta$ = recovery rate.

Three recovery rate definitions exist:

\vspace{0.3em}
\noindent\textbf{Recovery of Face Value (RFV):}

\textbf{Definition:} Fraction of defaultable claim nominal at default time.

\textbf{Formula:} $D(T,T) = 1_{\tau>T} + \frac{\delta}{P(t,T)} 1_{\tau \leq T}$ \quad (8.4)

\textbf{Usage:} Rating agencies (e.g., Moody's) for historical recovery rates.

\textbf{Pro:} Eases comparison across claims; also accommodate the delay required to resolve the default event.

\textbf{Con:} Makes $\delta$ sensitive to timing of default.

\vspace{0.3em}
\noindent\textbf{Recovery of Treasury Value (RTV):}

\textbf{Definition:} Fraction of equivalent risk-free zero-coupon bond at default time.  In case of default, the claimholder receives $\delta P(\tau, T)$ at time $\tau$ instead of 1 at time $T$.

\textbf{Formula:} $D(T,T) = 1_{\tau>T} + \delta 1_{\tau \leq T}$ \quad (8.5)

\textbf{Pro:} Makes $\delta$ insensitive to timing of default; simplifies valuation of defaultable claims.

\textbf{Con:} Not how recovery is determined in practice.

\vspace{0.3em}
\noindent\textbf{Recovery of Market Value (RMV):}

\textbf{Definition:} Fraction of pre-default bond value at default time. $\delta$ = proportional value loss incurred immediately by claimholder.

\textbf{Formula:} $D(T,T) = 1_{\tau>T} + \delta \frac{D(\tau^-, T)}{P(t,T)} 1_{\tau \leq T}$ \quad (8.6)

\textbf{Theory (Duffie-Singleton 1999):} Under reduced-form approach, RMV allows expressing defaultable claim value as EMM expectation of its terminal payoff discounted at risk-adjusted rate $R_t = r_t + s_t$ (where $s_t$ = credit spread):

$$D(t,T) = \mathbb{E}^Q\left[\exp\left(-\int_t^T (r_u + s_u)\,du\right)\right] \quad \text{(8.7)}$$

\textbf{Pricing implication:} Valuation of defaultable claims uses same tools as yield curve and zero-coupon bond pricing (factor models or HJM framework) for dynamics of $\{r_t\}$ and $\{s_t\}$.

\vspace{0.3em}
\noindent\textbf{Stylized Facts on Recovery Rates:}

\textbf{Priority impact:}

\noindent\textit{Bank loans:} First lien $>$ Senior unsecured $>$ Second lien

\noindent\textit{Bonds:} First lien $>$ Second lien $>$ Senior unsecured $>$  Senior subordinated$\approx$ Subordinated  $>$ Junior subordinated

\noindent Bank loans recover more than bonds at equivalent priority levels in the capital structure.

\textbf{Negative correlation:} There is a negative correlation between default likelihood and recovery. Or to say positive correlation between default likelihood and loss given default. 

\textbf{Amplification effect:} since Credit spread = default probability $\times$ loss given default,  the positive correlation between these two factors amplifies 
the variations in credit spreads.


\subsection{Corporate Bonds and Loans}
\subsubsection{The Debt Contract}

Corporations can borrow funds by issuing bonds on a market or by subscribing loans from a financial institution. The bond indenture or the loan structure define the respective rights and obligations of the borrower and the lender. The standard provisions are the nominal amount, the maturity, and the coupon payment schedule. Beyond the standard provisions, corporate bond and loan contracts are often augmented with covenants.

When a company deals with several types of creditors (through different debt issues), a usual contract feature is the \textbf{pari passu clause}, which stipulates that all contractual provisions should apply uniformly to all creditors (pari passu is a Latin expression for ``equal footing''). The pari passu clause does not mean that all creditors have the same rights. As a matter of fact, there exists a variety of provisions and covenants that creates substantial discrepancies across the cash flows that creditors are entitled to. The clause, however, implies that obligors must give to their existing creditors the same benefits that they would grant to future lenders. A consequence of the pari passu clause is that if an issuer wishes to revise the debt contract terms, she must do it with the knowledge of all creditors.

Thus, the pari passu clause deters shareholders and bondholders from acting strategically at the expense of one another. On one hand, it prevents shareholders from conducting a bondholder wealth expropriation by renegotiating debt terms in a bilateral manner, taking advantage of a lack of coordination among creditors. On the other hand, it prevents bondholders from engaging into a creditors' race by accelerating the liquidation of assets at the first signs of financial distress. Overall, the clause encourages collective bargaining with all creditors during debt restructurings (whether through private renegotiations or through the bankruptcy procedure).

Another common clause is the \textbf{negative pledge covenant}, which prevents the borrower from pledging assets to another lender. By extension, this clause forbids the issuance of new debt claims with equal or higher priority to existing debt.


\textbf{Covenants are classified into four categories:}
\begin{itemize}
    \item \textbf{Restrictions on investments:} These covenants include the prohibition of certain investment activities, a control on mergers and acquisitions, as well as restrictions on asset sales (asset sweep covenant). In addition, the net-worth covenant (a.k.a. safety covenant) requires that the issuer maintains a minimal net worth ratio (or tangible net worth ratio). Other covenants set limits on financial ratios like the interest coverage ratio, the debt-to-equity ratio, the current ratio, or the EBITDA-debt ratio to name a few. Finally, collateralization (secured debt) is a covenant that reduces the free usage of pledged assets.
    \item \textbf{Restrictions on shareholders' cash flows:} These covenants put constraints on the dividend policy. They can also forbid any share repurchase operation.
    \item \textbf {Debt priority structure:} In addition to the negative pledge covenant, a given level of priority can be assigned to different categories of lenders. The most common ones are senior, senior subordinated, subordinated and junior subordinated. Note that secured creditors have priority over unsecured creditors.
    \item \textbf{Option-like provisions:} These covenants include the borrower's right to repay early (callable debt), the lender's right to demand early repayment (putable debt), and the lender's right to convert her debt claim into equity (convertible debt). Also belongs to this category the sinking fund provision, which allows the issuer to periodically amortize the repayment of principal (at a pre-specified sinking fund price). To do so, the company transfers payments to a trustee who repurchases portions of the existing bonds. Thus, the issuer saves on interest payments and reduces the default risk on the debt (creditors have a claim on the fund as a collateral).
\end{itemize}

\subsubsection{Syndicated Loans}

Loan syndication is a \textbf{hybrid form} of debt placement between loans and bonds. A debt agreement is \textbf{initiated by the borrower and a bank (the arranger)} and marketed through a syndicate. The decision to syndicate the loan may occur before or after the deal is closed. In any case, \textbf{the terms of the deal (face value, maturity, covenants, pricing...)} are solely negotiated by the borrower and the arranger.

Then this arranger uses its network to distribute shares of the loan to banks that agree to become \textbf{members of the syndicate} (linked by pari passu clauses). The search for members (and therefore the construction of the syndicate) on the part of the arranger can be done \textbf{with a firm} \textbf{commitment} or on a \textbf{best effort basis}. In the first case, the arranger agrees to take on the undistributed shares of the loan. The borrower is therefore certain to raise all the funds he wants. The arranger will charge higher fees for this service.

\textit{Syndicate structure:}

\begin{itemize}
\item \textbf{The leader} (a.k.a. arranger, lead manager...)
\item \textbf{Co-agents} are responsible for an administrative task in the syndicate. Among those roles are: Documentation agent (drafts documents), administrative agent (calculates interests and collects repayment), collateral agent.
\item \textbf{Members.}
\end{itemize}






\subsection{Credit Ratings}
The senior unsecured bond is primarily evaluated, and the ratings of other issues are 
adjusted by notching for consistency. It is  based on accounting data, market data, historical default rates, and judgmental factors.

Moody's uses a rating system combining letters (Aaa, Aa, A, Baa, Ba, B, Caa, Ca, C) with numerical modifiers (1, 2, 3) following the letters. S\&P uses a system combining letters (AAA, AA, A, BBB, BB, B, CCC, CC, C) with signs (+, -, or null) following the letters.


There are \textbf{three types of credit events} recorded by rating agencies.

\begin{itemize}
\item \textbf{A missed or delayed payment} of interest and/or principal.

This credit event is sometimes referred to as a \textbf{technical default}. Though a missed payment is essentially a breach of the debt contract, its consequences can vary substantially depending on its interpretation. The market might consider that it only reflects a temporary liquidity shortage, or it can view it as an early sign of a more serious default.

\item \textbf{A distressed exchange}

This is the \textbf{outcome of renegotiations through a private workout}, and it is commonly referred to as a \textbf{debt restructuring}. Creditors agree to exchange their original claim with a modified one. A distressed exchange can be a \textbf{debt-for-debt swap}, in which case the new debt claim imposes lighter obligations (e.g., reduced coupon, extended maturity, foregone fraction of principal...). It can also be a \textbf{debt-for-equity swap}, in which case \textbf{lenders fully waive their creditors' rights and become owners of an agreed-upon fraction of the capital}.

In practice, \textbf{distressed exchanges can be difficult to implement when there are many creditors because of the holdout problem}.

\textbf{Holdout problem:} A form of freeriding behaviour. When there are many creditors, collectively all will be better off negotiating and easing some terms so the company would not default. But individually each of them are better off keeping their own previous deal and putting the burden on other creditors.

\textbf{The solution:} introduce fundamental rights : nominal, coupon and  maturity and all else are non fundamental (provisions and covenants ). And say that no one can force a creditor to change the fundamental ones. The non fundamental ones can be used as leverage by those saying yes to convince the one saying no!

\item \textbf{A bankruptcy filing}

This is when the distressed company officially announces that it files under the bankruptcy protection to begin a Court-supervised renegotiation process.
\end{itemize}
Credit ratings are not only estimating the default likelihood. They also take the expected loss given 
default into account. They are typically calculated as “through the business cycle” measures. That 
is, they are meant to convey an indication about creditworthiness for an investor wishing to hold 
the bond over a mid-run horizon. Although ratings aim at reflecting a forward-looking credit 
assessment, their revisions are not as frequent as market price updates. To incorporate more 
timely information, rating agencies also issue “watch lists”, signaling their analysis effort is focused 
on some issues whose rating might soon be corrected. historically lower ratings have been associated with likelihood of default and lower recovery rates. 


\subsection{Credit Derivatives }
\subsubsection{CDS Contract (credit default swap) }

With a \textbf{single-name CDS}, a \textbf{protection-buyer pays a periodic (usually quarterly) premium} ($p_{CDS}$) to a \textbf{protection-seller} who, in turn, compensates for the loss in case of a default event affecting an issuer.

The CDS is entirely characterized by a \textbf{maturity}, a \textbf{notional amount}, a \textbf{reference entity} (the issuer of defaultable securities), reference obligations, and a list of credit events. Settlement of CDS can be carried through \textbf{physical settlement}, but this might be hampered by the limited number of outstanding bonds. \textbf{Cash settlement} may also be difficult to implement because of the lack of liquidity for defaulted bonds. The third possibility is a \textbf{settlement auction}.

Before the \textbf{2007-2008 financial crisis}, the premium of the CDS was quoted so that it corresponds to a \textbf{zero value for the CDS upon inception}. This market convention has been amended since, and entering a CDS contract now requires the payment of an upfront payment.


An \textbf{index CDS} is a contract \textbf{written on a credit index}. The \textbf{buyer of the index CDS receives a periodic premium} (for investing in a credit portfolio). The \textbf{seller of the CDS} is compensated for the default loss for each credit event affecting the set of references underlying the credit index.

The two major credit index families are the \textbf{Dow Jones CDX in North America} and the \textbf{iTraxx in Europe}. Regarding CDX indices, the \textbf{CDX North America investment grade (CDX NA IG)} is a portfolio of 125 single-name CDS. The \textbf{CDX North America high yield} is a portfolio of 100 single-name CDS. the iTraxx Europe is a portfolio of 125 single-name CDS. The iTraxx Crossover contains 50 high-yield references.

\subsection{Quoting Convention}

Instead of quoting a CDS premium which varies from one risky underlying reference to another, CDS contracts now have an annual fixed coupon of either 100 bps for investment grade or 500 bps for high yield/speculative grade. Let $\tau$ denote the payment frequency (usually the quarter) and the $t_j$, $j = 1, \ldots, n$ correspond to the CDS payment dates (contract maturity is $T = n\tau$).

The evaluation of the CDS contract requires the specification of a \textbf{cumulative survival function} $CS(t)$ that represents the \textbf{probability that the reference entity does} \textbf{not default until time} $t$. In the absence of arbitrage, this function is calculated under the Equivalent Martingale Measure. The ISDA model posits a simple \textbf{exponential decay}, i.e., $CS(t) = \exp(-\lambda t)$, where $\lambda$ is referred to as the \textbf{hazard rate} and represents the \textbf{instantaneous default probability (under the EMM)}.

According to the ISDA model, the present value of the fixed leg of the CDS can be written as:

$$V_1 = p_{CDS} \sum_{j=1}^{n} \tau CS(t_j)d(t_j) + \frac{1}{2}p_{CDS} \sum_{j=1}^{n} \tau [CS(t_{j-1}) - CS(t_j)]d(t_j), \quad (8.8)$$

where $d(t_j)$ is the discount factor up to time $t_j$. The second term of equation (8.8) accounts for the \textbf{accrued CDS premium that must be paid in scenarios where default occurs between two} \textbf{payment dates}. The ISDA model proxies the accrued premium by half of $p_{CDS}$.

Using a similar logic, the present value of the variable leg of the CDS is:

$$V_2 = (1-R) \sum_{j=1}^{n} [CS(t_{j-1}) - CS(t_j)]d(t_j), \quad (8.9)$$

where $R$ stands for the recovery rate.

The fair pricing of the CDS entails that $V_1 = V_2$. Let $\lambda^*$ denote the hazard rate resulting from the CDS fair price. The corresponding cumulative survival function is denoted by $CS^*(t)$.

Turning to the new quoting convention, the fair pricing of the CDS contract with fixed coupon $c_{CDS}$ (equal to 100 bps for investment grade or 500 bps for high yield/speculative grade) and upfront payment $U$ entails that:

$$U = -c_{CDS} \sum_{j=1}^{n} \tau CS^*(t_j)d(t_j) - \frac{1}{2} c_{CDS} \sum_{j=1}^{n} \tau [CS^*(t_{j-1}) - CS^*(t_j)]d(t_j)$$
$$+ (1 - R) \sum_{j=1}^{n} [CS^*(t_{j-1}) - CS^*(t_j)]d(t_j). \quad (8.10)$$

Thus, using equations (8.9) and (8.10):

$$U = (p_{CDS} - c_{CDS}) \sum_{j=1}^{n} \tau CS^*(t_j)d(t_j)$$
$$+ \frac{1}{2}(p_{CDS} - c_{CDS}) \sum_{j=1}^{n} \tau [CS^*(t_{j-1}) - CS^*(t_j)]d(t_j). \quad (8.11)$$

The upfront payment compensates for substituting the CDS premium with the fixed coupon. This substitution induces a difference in present values calculated from the fair pricing of the CDS. Note that the upfront payment is positive when $p_{CDS} > c_{CDS}$. But it can also be negative, i.e., the buyer of the CDS receives a positive cash flow, if the fixed coupon overcompensates for the credit risk exposure.

Note that by using the product of discount factors with cumulative survival functions, the ISDA model implicitly assumes the independence between interest rate risk and credit risk.



\subsubsection{CDS-bond Basis}

Since credit risk can be traded on the corporate \textbf{bond market} and on the \textbf{CDS market}, the credit spreads implied from the two markets should be quite the same to bar arbitrage. Consider the following strategy:

\begin{itemize}
\item \textbf{Buy the CDS protection},
\item \textbf{Buy a floating rate corporate bond},
\item \textbf{Sell the bond on the repo market.}
\end{itemize}

In the absence of arbitrage, we therefore must have that:

$$p_{CDS} = y - r. \quad (8.12)$$

The \textbf{CDS-bond basis} is defined as the difference between the CDS spread and the corporate bond yield spread, and it should theoretically be nil in the absence of arbitrage.

There are, however, several \textbf{practical reasons for the CDS-bond basis to} \textbf{deviate from zero}. First, corporations seldom issue floating rate bonds with maturities that coincide exactly with those of CDS contracts. The corporate bond yield spread must therefore be recovered from fixed rate bond issues with specific maturities, and this causes some measurement problems to estimate $y$. Second and more importantly, \textbf{there are limits to arbitrage}, which prevent market participants to constantly implement the strategy described above. These frictions include trading liquidity, funding cost, counterparty risk, and collateral quality.

In practice, the CDS-bond basis for a reference $j$ and maturity $T$ is measured as follows:

$$\text{Basis}_j(T) = p_{CDS,j}(T) - PECDS_j(T), \quad (8.13)$$

where $PECDS_j(T)$, the Par Equivalent CDS, \textbf{is the theoretical bond-implied CDS spread}.

The method to compute $PECDS_j(T)$ is related to the ISDA model. First, one \textbf{determines the} \textbf{cumulative survival function} $CS^*(t)$ ensuring the fair pricing of the CDS (i.e., equalizing the legs equations (8.8) and (8.9)). Then, one solves for the smallest shift in the term structure of $CS^*(t)$ that reconciles the theoretical bond price with its market value. That is, the bond-implied cumulative survival function is:

$$CS_{\text{bond}}(t) = CS^*(t) + \varepsilon, \quad (8.14)$$

with

$$\varepsilon = \text{argmin } [PV(CS_{\text{bond}}(t)) - \text{Market price}]^2, \quad (8.15)$$

Finally, $PECDS_j(T)$ is calculated as the par equivalent CDS spread that uses $CS_{\text{bond}}(t)$ as the cumulative survival function. That is, combining equations (8.8) and (8.9):

$$PECDS_j(T) = \frac{(1 - R) \sum_{j=1}^{n}[CS_{\text{bond}}(t_{j-1}) - CS_{\text{bond}}(t_j)]d(t_j)}{\sum_{j=1}^{n} \tau CS_{\text{bond}}(t_j)d(t_j) + \frac{1}{2}\sum_{j=1}^{n} [CS_{\text{bond}}(t_{j-1}) - CS_{\text{bond}}(t_j)]d(t_j)} \quad (8.16)$$

\textbf{The CDS-bond basis has typically been negative on the U.S. corporate bond market.}






\subsection{CLO (Collateralized Loan Obligations)}

The \textbf{CLO is an investment structure pooling a collection of defaultable assets} (loans, bonds, mortgage-backed securities, etc.) via a special purpose vehicle. \textbf{All the assets are securitized and sold by tranches}. Typical tranches are:

\begin{itemize}
\item The \textbf{senior tranche}, gathering around 70\% to 90\% of debt nominal,
\item \textbf{Other debt tranches} (known as mezzanine), gathering around 5\% to 15\% of debt nominal,
\item \textbf{The equity (subordinate) tranche}, gathering around 2\% to 15\% of debt nominal.
\end{itemize}

The securitization of CLOs allows banks to market illiquid loans and increase their lending activity. But it may also provide them an incentive to get rid of bad quality loans (\textbf{moral hazard issue}). The different tranches are designed to attract various investor profiles with different risk/return objectives.

Over the maturity period of the CLO, \textbf{defaults on the pool of underlying loans will impact on the payoff of tranches depending on their seniority level}. This, the first default will reduce the payoff to the equity tranche. Then, subsequent defaults will gradually affect the payoff to mezzanine and ultimately senior tranches.

As a consequence, while the \textbf{risk borne by the equity tranche is essentially related to the highest level of default risk among the underlying loans}, the \textbf{risk borne by senior tranches is mostly related to default clusters and is a joint effect of total default risk and default correlation risk}.




\section{Structural Models}
Structural models of credit risk posit the dynamics of a fundamental state variable (typically the 
asset value or a cash flow of the issuing firm) and resort to contingent claims analysis to evaluate 
defaultable claims.
 

In all cases, the \textbf{preference-free pricing rule} holds if we can replicate the payoff (and assuming no arbitrage). even if the state variable (the firm's assets or cash flows) are not traded, \textbf{provided that one contingent claim (typically the stock) is traded } we can do that. Also if the value of assets is strongly correlated with a market variable (something tradable), like a mining firm or energy firm, the assumption holds. The structural approach can therefore be applied to all securities issued by \textbf{publicly traded firms}.

\subsection{The Merton Model and its Extensions}

\subsubsection{Zero-coupon Debt}

\textbf{Merton's (1974) model}: Consider a firm with a very simple capital structure: equity and a zero-coupon debt with nominal amount $M$ and maturity $T$. Let $\{v_t\}_{t \geq 0}$ denote the value process of the firm's \textbf{assets}. The \textbf{terminal wealth of shareholders and debtholders is:}

$$S_T = \max(v_T - M, 0), \quad (9.1)$$

$$D_T = \min(v_T, M) = M - \max(M - v_T, 0). \quad (9.2)$$

Using analogies with options, \textbf{equity is a call option on the value of assets and debt provides a position on the nominal amount minus a put option on assets}. By taking up the \textbf{assumptions of Black and Scholes (1973) and assuming in particular that the value of the firm's assets is governed by a geometric Brownian motion with volatility} $\sigma$, we obtain:

$$S = v\Phi(d_1) - Me^{-rT}\Phi(d_2), \quad (9.3)$$

$$D = Me^{-rT}\Phi(d_2) + v\Phi(-d_1), \quad (9.4)$$


with

$$d_1 = \frac{\ln\left(\frac{v}{M}\right) + \left(r + \frac{\sigma^2}{2}\right)T}{\sigma\sqrt{T}}, \quad d_2 = \frac{\ln\left(\frac{v}{M}\right) + \left(r - \frac{\sigma^2}{2}\right)T}{\sigma\sqrt{T}}, \quad (9.5)$$

($\Phi(d_2)$ is the EMM probability of no default and $\Phi(-d_2) = 1 - \Phi(d_2)$ is the EMM probability of default),

\textbf{The yield spread (assimilated to the credit spread) is given by:}

$$s(T) = -\frac{1}{T}\ln\left(\frac{D}{M}\right) - r, \quad (9.6)$$

The term structure of credit spreads created by this model can be \textbf{increasing} (for low debt levels), \textbf{decreasing} (for high debt levels), or \textbf{bell-shaped} (for intermediate debt levels). However, the TSCS is \textbf{constrained to converge to zero at the short end}. This property holds because default is predictable, i.e., as $T \to 0$ the probability of immediate default approaches 0 (if $v_0 \geq M$) or 1 (if $v_0 < M$). This is not consistent with empirical evidence, and can be corrected by introducing a jump component in the dynamics.

\textbf{Merton (1976) Jump Model-Asset value dynamics (under EMM):}
$$\frac{dv_t}{v_t} = (r - \lambda\mu_J)dt + \sigma dW_t + (e^J - 1)dN_t \quad (9.7)$$

where $\{W_t\}$ is a standard Brownian motion, $\{N_t\}$ is a Poisson process with intensity $\lambda$, and $J$ is the log-jump magnitude with $\mathbb{E}[e^J - 1] = \mu_J$.

Equity is still viewed as a European call on firm assets with strike $M$ and maturity $T$. The European call premium is:
$$S = \sum_{n=0}^{\infty} \frac{(a(1+\mu_J)T)^n}{n!}e^{-\lambda(1+\mu_J)T}\text{Call}_{BS}\left(v, T, M, r -\lambda \mu_J + \frac{n\ln(1+\mu_J)}{T}, \sigma^2 + \frac{n\sigma^2}{T}\right) \quad (9.8)$$

where $\text{Call}_{BS}$ is the Black-Scholes call premium  with modified risk-free rate and volatility, each call being weighted by the 
probability of observing n jumps within T. The infinite sum can be truncated at a relatively low level (e.g., first 50 terms) for sufficiently accurate evaluation.

\textbf{Balance sheet equality:} $D = v - S$ (debt = assets - equity).

\subsubsection{Secured Debt}

\textbf{Collateral effect:} Adding collateral (guarantee or pledge) reduces risk exposure of creditors since in event of default they recover the value of the asset put up as collateral.

\textbf{Constant collateral value:} Assume collateral value is constant, denoted by $K$ (with $K < M$). Terminal value of secured debt:
$$D_T^a = M \times \mathbf{1}_{v_T \geq M} + K \times \mathbf{1}_{v_T < M}$$

Creditors recover proceeds of the sale of collateral as priority and nothing else. If $K$ does not cover entire claim, remaining debt is assigned lowest priority rank among claims. Secured debt involves position in cash-or-nothing call and put. Under Black-Scholes assumptions:
$$D^a = Me^{-rT}\Phi(d_2) + K\Phi(-d_2) = e^{-rT}[K + (M-K)\Phi(d_2)]$$

Adding collateral generally reduces risk exposure but not necessarily. Amount covered by collateral is guaranteed but remaining debt becomes very risky (zero recovery). Net effect on debt risk depends on recovery obtained with collateral versus recovery without collateral given priority ranking.

\textbf{Stochastic collateral value:} Collateral is not always constant. When pledged asset is financial security or raw material, value is random. Assume collateral is traded asset with price process $\{K_t\}$ following geometric Brownian motion, correlated with asset value:
$$\frac{dK_t}{K_t} = \mu_K dt + \eta dZ_t, \quad dZ_t dW_t = \rho dt$$

Terminal value of secured debt:
$$D_T^a = M \times \mathbf{1}_{v_T \geq M} + K_T \times \mathbf{1}_{v_T < M}$$

Current value under Black-Scholes:
$$D^a = Me^{-rT}\Phi(d_2) + K\Phi(-d_2 - \eta\sqrt{T})$$

\textbf{Wrong Way Risk:} Correlation $\rho$ between collateral and assets plays critical role in risk reduction effectiveness. If $\rho$ is very high, collateral has little value when assets depreciate, so it will not effectively protect against default risk.

\subsubsection{Interim Payments}

\textbf{Setup (Geske 1977):} Consider a debt of maturity $T$ generating payments $c_i$ on dates $t_i$, $i = 1, \ldots, n$ (with $T = t_n$). Notations $v_t$, $S_t$, and $D_t^c$ designate values of assets, equity and debt with coupons on date $t_i$.

\textbf{Equity as compound call:} Equity is interpreted as a call option $n$ times compounded, written on the assets. Generalizing Merton's formula, the value of the coupon bond is written as:
$$D^c = \sum_{i=1}^{n} c_i e^{-rt_i}\Phi_i(h_i, \{\theta_{ij}\}) + v\left(1 - \Phi_n\left(h_i + \sigma\sqrt{t_i}, \{\theta_{ij}\}\right)\right) \quad \text{(9.15)}$$

with, for all $i = 1, \ldots n$ and all $j = 1, \ldots n$:
$$h_i = \frac{1}{\sigma\sqrt{t_i}} \left[\ln\frac{v}{c_i} + \left(r - \frac{\sigma^2}{2}\right)t_i\right] \quad \text{(9.16)}$$

$$\theta_{ij} = \frac{\sqrt{t_i}}{\sqrt{t_j}}, \quad i < j \quad \text{(9.17)}$$

$\Phi_i(x_i, \{y\})$ denotes the cumulative distribution function of the standard multivariate normal law (of dimension $i$), evaluated at points $x_i$, with a variance-covariance matrix $y$.

\textbf{Critical default thresholds:} The terms $v_i^*$ represent the critical thresholds for the value of the assets at which the company defaults. They are determined according to the method of financing the debt service. In the case where the repayment of the debt cannot be financed by selling assets, then the shareholders will continue to pay the intermediate coupons if they can raise the necessary funds by issuing shares. The characterization of the critical default threshold is then given by:
$$S(v_i^*) = c_i \quad \text{(9.18)}$$

The computation of debt value is solved recursively. At date $T$, the default threshold is $c_n$. From the expression of $D_{n-1}$, we obtain that of $S_{n-1} = v_T - D_{n-1}$. Then numerically solve the equation $S_{n-1} = c_{n-1}$ to infer the critical threshold $v_{n-1}^*$. We then continue with the expression of $D_{n-2}$ and that of $S_{n-2}$. Solving the equation $S_{n-2} = c_{n-2}$ knowing $v_{n-1}^*$ allows us to extract $v_{n-2}^*$, and so on until the initial date. Solving the Geske (1977) model is highly computational and can become complex if the number of intermediate payments is high. Approximation methods can be useful.

\textbf{Approximation by zero-coupon decomposition:} A simplifying approach is to treat coupon debt as the sum of zero-coupon debts, the latter being valued directly using Merton's formula. Thus:
$$D^c = \sum_{i=1}^{n} D_i \quad \text{(9.19)}$$

This method is approximate in that it considers the default events on the different payments as independent. However, in practice, the payment of the $(k + 1)$th coupon is conditional on the fact that the $k$ previous coupons have been paid. The numerical example below shows that this approximation gives an acceptable result when the number of coupons is modest.

$$D_i = c_i e^{-rt_i}\Phi(d_{2,i}) + v\Phi(-d_{1,i}) \quad \text{(9.20)}$$

$$d_{1,i} = \frac{\ln(v/c_i) + (r + \frac{\sigma^2}{2})t_i}{\sigma\sqrt{t_i}}, \quad d_{2,i} = \frac{\ln(v/c_i) + (r - \frac{\sigma^2}{2})t_i}{\sigma\sqrt{t_i}} \quad \text{(9.21)}$$

\textbf{Application to Short- and Long-term Debt:}

Consider a company with two debt issues, one with nominal value $m$ and maturity $t$ (short-term debt), the other with nominal value $M$ and maturity $T$ (long-term debt). The two debts have the same priority. In the event of a short-term default, the creditors will share a fraction $\theta$ of the value of the assets in proportion to their nominal contributions. The terminal value of the short-term debt is:

$$d_t = m \times \mathbf{1}_{v_t \geq m} + \frac{m}{m+M} \theta v_t \times \mathbf{1}_{v_t < m} \quad \text{(9.22)}$$

$$D_T = M \times \mathbf{1}_{v_T \geq m} \times \mathbf{1}_{v_T \geq M} + \theta v_T \times \mathbf{1}_{v_t \geq m} \times \mathbf{1}_{v_T < M} + \frac{M}{m+M} \theta v_t \times \mathbf{1}_{v_t < m} \quad \text{(9.23)}$$

We therefore obtain the following formulas:

$$d = me^{-rt}\Phi(z_2) + \frac{m}{m+M} \theta \nu \Phi(-z_1), \quad \text{(9.24)}$$

$$D = Me^{-rT}\Phi_2\left(z_2, z_4; \frac{\sqrt{t}}{\sqrt{T}}\right) + \theta \nu \Phi_2\left(z_1, -z_3; -\frac{\sqrt{t}}{\sqrt{T}}\right) + \frac{M}{m+M} \theta \nu \Phi(-z_1), \quad \text{(9.25)}$$

with

$$z_1 = \frac{\ln(v/m) + (r + \frac{\sigma^2}{2})t}{\sigma\sqrt{t}}, \quad z_2 = \frac{\ln(v/m) + (r - \frac{\sigma^2}{2})t}{\sigma\sqrt{t}} \quad \text{(9.26)}$$

$$z_3 = \frac{\ln(v/M) + (r + \frac{\sigma^2}{2})T}{\sigma\sqrt{T}}, \quad z_4 = \frac{\ln(v/M) + (r - \frac{\sigma^2}{2})T}{\sigma\sqrt{T}} \quad \text{(9.27)}$$

\subsubsection{Net Worth Covenant}

It requires the company to maintain its net worth above a certain level during the life of the debt contract. Net worth generally refers to the book value of the assets (from which the value of intangible assets is sometimes subtracted) less all liabilities. When debts remain stable, this clause is equivalent to setting a minimum threshold for the value of the assets. 
Let $H$ denote the threshold above which the asset value process must remain. If $v_t$ crosses the threshold $H$ before debt maturity $T$, the company is in default. In optional terms, equity is a call option written on the assets with a knock-out barrier $H$ (down and out call).

Under the Black-Scholes assumptions for asset dynamics and the Merton (1974) model for the capital structure of the firm, the value of equity is given by:

$$S = v\Phi(d_1) - Me^{-rT}\Phi(d_2) - v\left(\frac{H}{v}\right)^{2r/\sigma^2 + 1}\Phi(d_3) + Me^{-rT}\left(\frac{H}{v}\right)^{2r/\sigma^2 - 1}\Phi(d_4), \quad \text{(9.28)}$$

where the terms $d_1$ and $d_2$ keep their same definition and

$$d_3 = \frac{\ln(H^2/(vM)) + (r + \frac{\sigma^2}{2})T}{\sigma\sqrt{T}}, \quad d_4 = \frac{\ln(H^2/(vM)) + (r - \frac{\sigma^2}{2})T}{\sigma\sqrt{T}} \quad \text{(9.29)}$$

The debt value is obtained by balance sheet equality $D = v - S$. From the creditors' point of view, the minimum net worth clause reduces the risk exposure because it precipitates the default of the company.


\subsubsection{Debt Priority }
Determines the creditor's rank in the liquidation procedure. More generally, it influences the relative importance of each creditor during renegotiations in the event of financial distress. Senior and junior debts represent the two main ranks of priority.

Secured debt has the highest rank of priority. In the event of liquidation, the collateral asset is first 
sold. If the proceeds from this sale do not cover the nominal amount of the debt, the remaining 
claim joins those of the other lenders. Otherwise, the excess amount of money is distributed to 
the other creditors according to their rank of priority. secured debt very often contains a negative pledge covenant. This 
clause prohibits the company from putting an asset as collateral for a new debt.

\textbf{Senior and Junior Debt Formulas:} Consider a company whose capital structure consists of equity, a senior debt of nominal $M_1$ and maturity $T$ and a junior debt of nominal $M_2$ and same maturity $T$. The shareholders' decision to default is based on the total debt $M_1 + M_2$ (pari passu clause). Consequently, the terminal value of equity is given by:

$$S_T = \max(v_T - (M_1 + M_2), 0). \quad \text{(9.30)}$$

The terminal value of senior debt is:

$$D_T = \min(M_1, v_T) = M_1 - \max(M_1 - v_T, 0), \quad \text{(9.31)}$$

and from balance sheet equality the terminal value of junior debt is:

$$d_T = v_T - S_T - \min(M_1, v_T) = \max(v_T - M_1, 0) - \max(v_T - (M_1 + M_2), 0). \quad \text{(9.32)}$$

The current values of senior and junior debt are given by:

$$D = M_1 e^{-rT}\Phi(a_2) + v\Phi(-a_1), \quad \text{(9.33)}$$

$$d = v\Phi(a_1) - M_1 e^{-rT}\Phi(a_2) - v\Phi(b_1) + (M_1 + M_2)e^{-rT}\Phi(b_2), \quad \text{(9.34)}$$

with

$$a_1 = \frac{\ln \frac{v}{M_1} + (r + \frac{\sigma^2}{2})T}{\sigma\sqrt{T}}, \quad a_2 = \frac{\ln \frac{v}{M_1} + (r - \frac{\sigma^2}{2})T}{\sigma\sqrt{T}}, \quad \text{(9.35)}$$

$$b_1 = \frac{\ln \frac{v}{M_1 + M_2} + (r + \frac{\sigma^2}{2})T}{\sigma\sqrt{T}}, \quad b_2 = \frac{\ln \frac{v}{M_1 + M_2} + (r - \frac{\sigma^2}{2})T}{\sigma\sqrt{T}}. \quad \text{(9.36)}$$


\subsubsection{ Stochastic Interest Rates}
Consider the Merton model with a Heath-Jarrow-Morton yield curve. dynamics under the EMM:

$$\frac{dv_t}{v_t} = r_t dt + \sigma dW_t, \quad \text{(9.37)}$$

$$\frac{dP(t,T)}{P(t,T)} = r_t dt - \sigma_p(t,T) dZ_t, \quad \text{(9.38)}$$

with $\sigma_p(t,T)$ a deterministic function of time and $dW_t dZ_t = \rho dt$.

$$S = v\Phi(z_1) - MB(0,T)\Phi(z_2), \quad \text{(9.39)}$$

with

$$z_1 = \frac{\ln \frac{v}{MB(0,T)} + \frac{1}{2}V(T)}{\sqrt{V(T)}}, \quad z_2 = z_1 - \sqrt{V(T)}, \quad \text{(9.40)}$$

$$V(T) = \int_0^T (\sigma^2 + \sigma_p^2(u,T) + 2\rho\sigma\sigma_p(u,T)) du. \quad \text{(9.41)}$$

The stochastic nature of interest rate hardly affects the credit spreads for short horizons and the 
impact on long-term credit spreads is limited. By contrast, the 
correlation between the bond returns and asset returns affects almost the whole term structure 
and the effect is much stronger 


\subsection{Option-like Provisions }
\subsubsection{Callable Debt }
shareholders have the right (by design, of American type) to redeem the debt at any time before maturity for a price determined in advance.
The call option allows the borrowing company to better manage interest rate risk (redeemming the debt when interest rates are low and to refinance its debt at a lower cost). Also In the event of a conflict with lenders, shareholders can simply buy back the debt  and free themselves from the constraints imposed by the contract.

\textbf{CRR Approach:} debt of nominal amount $M$, maturity $T$ redeemable at any time by shareholders. In general, the redemption price can be any deterministic function of time, we will simply designate it by $K$ (at maturity, we must have $K = M$).

Let $j$ $(0 \leq j \leq n)$ denote the position in the tree, numbered from the bottom (position 0) to the top (position $n$). we number time backwards, from $k = 0$ (maturity $T$) to $k = n$ (the initial time):

$$G_j^k = (\pi G_{j+1}^{k-1} + (1 - \pi)G_j^{k-1})e^{-r\frac{T}{n}} \quad \text{(9.42)}$$

where $\pi$ denotes the EMM probability of an upward movement in the underlying asset.

The valuation algorithm can be initialized at maturity using the terminal values of the callable debt or to ensure much better convergence, start the valuation at the penultimate date of the tree. At date $k = 1$, there will be no possibility of buying back the debt before maturity. The continuation value is therefore given by the Merton formula (1974) with a remaining maturity of $T/n$. Therefore, the algorithm starts with:

$$D_j^1 = \min\left(K, Me^{-r\frac{T}{n}}\Phi(d_2) + v_j^1\Phi(-d_1)\right), \quad \text{(9.43)}$$

$$v_j^1 = vu^{n-j-1}, \quad \text{(9.44)}$$

and $d_1$ and $d_2$ defined with $v_j^1$. For all other periods before maturity, the value of the callable debt is obtained by taking the minimum between its continuation value and the redemption price prevailing at that period. Thus, we have, for all $k \leq n$:

$$D_j^k = \min\left(K, (\pi D_{j+1}^{k-1} + (1-\pi)D_j^{k-1})e^{-r\frac{T}{n}}\right), \quad \text{(9.45)}$$

\textbf{CRR Tree Reminder:} time to maturity $T$ is divided into $n$ periods of equal duration $T/n$. In each period, the value of the firm's assets can either go upward by being multiplied by $u$ or downward by being multiplied by $d = \frac{1}{u}$. 

$$u = \exp(\sigma\sqrt{T/n}), \quad \pi = \frac{e^{r\frac{T}{n}} - d}{u - d} $$


\subsubsection{Convertible Debt }
Gives its holder the right to convert his security into a predetermined fraction of shares.
 The conversion of the debt gives 
rise to the issue of n new shares in addition to the N existing shares. The conversion ratio $q=\frac{n}{N+n}$ therefore measures the share of equity obtained by the creditors at the time of 
conversion. 

The terminal wealth of shareholders and creditors are:

$$S_T = \max(v_T - M, 0) - q \times \max\left(v_T - \frac{M}{q}, 0\right), \quad \text{(9.47)}$$

$$D_T = v_T - \max(v_T - M, 0) + q \times \max\left(v_T - \frac{M}{q}, 0\right). \quad \text{(9.48)}$$

The value of the European convertible debt is:

$$D^{\text{conv}}(v,P) = v\Phi(-d_1) + Me^{-rT}[\Phi(d_2) - \Phi(d_4)] + qv\Phi(d_3), \quad \text{(9.49)}$$

where $d_1$ and $d_2$ have their usual definition and:

$$d_3 = \frac{\ln \frac{qv}{M} + (r + \frac{\sigma^2}{2})T}{\sigma\sqrt{T}}, \quad d_4 = \frac{\ln \frac{qv}{M} + (r - \frac{\sigma^2}{2})T}{\sigma\sqrt{T}}. \quad \text{(9.50)}$$

We can use the closed-form expression for the European convertible debt to initiate the pricing algorithm for the American convertible debt using the Cox-Ross-Rubinstein binomial tree. More specifically:

$$D_j^1 = \max\left(qv_j^1, D^{\text{conv}}(v_j^1, T/n)\right), \quad \text{(9.51)}$$

with $v_j^1 = vu^{n-j-1}$. For all other periods before maturity, the recursive algorithm continues by comparing the continuation value with the exercise value. Note that we take the maximum in this case, because, when it comes to convertible debt, the conversion option is in the hands of the creditors. Thus, we have, for all $k \leq n$:

$$D_j^k = \max\left(qvu^{n-j-k}, (\pi D_{j+1}^{k-1} + (1-\pi)D_j^{k-1})e^{-r\frac{T}{n}}\right). \quad \text{(9.52)}$$


\subsubsection{Callable Convertible Debt}
The shareholders' decision to call the debt or not determines the terms under which the creditors 
will make their decision to convert or not. Subsequently, either the shareholders do not announce the repurchase of the debt, in 
which case the creditors will choose the maximum between the continuation value and the 
conversion value, or the shareholders announce the repurchase of the debt, in which case the 
creditors will choose the maximum between the repurchase value and the conversion value (the 
latter situation being referred to as forced conversion).

$$D_j^k = \min\left[\max\left(qvu^{n-j-k}, (\pi D_{j+1}^{k-1} + (1-\pi)D_j^{k-1})e^{-r\frac{T}{n}}\right), \max(qvu^{n-j-k}, K)\right]$$

The algorithm can be initialized in period $k = 1$ by using the Black-Scholes formula for convertible debt. Which yields:

$$D_j^1 = \min\left[\max\left(qv_j^1, D^{\text{conv}}(v_j^1, T/n)\right), \max(qv_j^1, K)\right]$$

\subsection{First Passage Time Models}

Default occurs the first time the value of the firm's assets falls below some specified threshold. Default time is then characterized by a random variable called a first passage time (or first hitting time).

$$\tau = \inf(t \geq 0: v_t = H_t)$$

When the process $v_t$ is driven by a geometric Brownian motion and the boundary is the constant $H$, the cdf of $\tau$ ($F_1(\tau) := Q(\tau < t)$) and the joint cdf of $\tau$ and $v$ ($F_2(\tau, v) := Q(\tau > t, v > x)$) are explicit:

$$F_1(\tau) = \Phi\left(\frac{\ln(H/v) - (r - \frac{\sigma^2}{2})t}{\sigma\sqrt{t}}\right) + \left(\frac{H}{v}\right)^{\frac{2r}{\sigma^2}-1}\Phi\left(\frac{\ln(H/v) + (r - \frac{\sigma^2}{2})t}{\sigma\sqrt{t}}\right), \quad \text{(9.56)}$$

$$F_2(\tau, v) = \Phi\left(\frac{\ln(v/x) + (r - \frac{\sigma^2}{2})t}{\sigma\sqrt{t}}\right) - \left(\frac{H}{v}\right)^{\frac{2r}{\sigma^2}-1}\Phi\left(\frac{\ln(H^2/(vx)) + (r - \frac{\sigma^2}{2})t}{\sigma\sqrt{t}}\right). \quad \text{(9.57)}$$

These results can readily be extended to the case of a deterministic exponential boundary of the form $H_t = K\exp(-\gamma(T - t))$. Indeed, using the solution of the geometric Brownian motion, we can write that: Therefore $\tau$ remains the first passage time of a geometric Brownian motion (with drift $r - \gamma$) to the constant boundary $K\exp(-\gamma T)$.

\subsubsection{Black-Cox Model }
Black and Cox (1976) solve for the pricing of the zero-coupon bond with nominal M, maturity T 
and exponential default boundary $H_t = K\exp(-\gamma(T - t))$.
The introduction of a default boundary alters default probabilities across all horizons. It therefore 
provides great flexibility in modelling the TSCS. Beyond the level of credit spreads, first passage 
time models generate a variety of shapes for the TSCS.

If the default barrier grows slowly relative to the risk-free rate, the TSCS is translated 
but essentially keeps the same shape. if the default boundary grows much faster than the risk-free rate, the shape of the TSCS 
can be distorted.

\subsubsection{Briys-de Varenne Model}

Briys and de Varenne (1997) evaluate zero-coupon bonds with an exponential default boundary when interest rates are stochastic in the Vasicek framework. The default boundary is directly linked to the Treasury discount bond:
$$v_b(t) = \lambda MP(t,T) \quad \text{for } t < T, \quad v_b(T) = M,$$

Under RTV assumption, upon default, creditors recover:
$$D_{bdv}(v_b) = \delta MP(t,T),$$
with $\delta < \lambda$ accounting for default costs.

The default boundary position greatly impacts the TSCS. Like the Black-Cox model, the effect of correlation between asset and rate dynamics matters as much (if not more) than interest rate volatility alone on the term structure of credit spreads.

\subsubsection{Collin-Dufresne-Goldstein Model }

In preceding models, leverage is static: since assets grow at rate $r$ under the EMM but debt is constant, leverage decreases over time, underestimating credit spreads. In practice, firms target leverage ratios, so debt burden should be adjusted dynamically.

Collin-Dufresne and Goldstein (2001) propose a dynamic default threshold:
$$dk_t = a(y_t - b - k_t)dt,$$
where $y_t = \ln(v_t)$ and $a, b$ are constants.

The default time is the first passage time of log-leverage $l_t = k_t - y_t$ to zero: $\tau = \inf(t \geq 0: l_t = 0)$.

The log-leverage dynamics rearrange to:
$$dl_t = a(L - l_t)dt - \sigma dW_t,$$
where $L = a^2/a - b$.

Thus default is the first passage time of a Vasicek process to zero. Under RTV (fraction $1-\omega$ recovered):
$$D_{cdg} = e^{-r\tau}(1 - \omega Q(\tau < T)).$$

\subsection{Estimation Methods}

We review three methods to estimate the current asset value and asset return volatility assuming assets are driven by a GBM using a time series of $m$ daily stock prices. Let $S$ denote the market capitalization (stock price $\times$ shares outstanding).

Equity is a call option on assets. Synthesizing the firm's liabilities to a single default point at the 1-year horizon, the amount due is proxied by short-term debt ($std$) plus half of long-term debt ($ltd$). Under Black-Scholes, market capitalization is:

$$S = v\Phi(d_1) - (std + 0.5ltd)e^{-r}\Phi(d_2) := f(v, \sigma), \quad \text{(9.79)}$$


$$d_1 = \frac{\ln\frac{v}{std + 0.5ltd} + r + \frac{\sigma^2}{2}}{\sigma}, \quad d_2 = \frac{\ln\frac{v}{std + 0.5ltd} + r - \frac{\sigma^2}{2}}{\sigma}$$

Setting $r$ equal to the 1-year risk-free rate.

\subsubsection{ System of Equations }
obtained from Itô's lemma: 

$$\sigma_S = \sigma \Phi(d_1) := g(v, \sigma). \quad \text{(9.80)}$$

Using observed $\sigma_S^{obs}$ and $\sigma_S^{obs}$ (historical or implied), one can solve the following system:

$$\begin{cases}
\sigma_S^{obs} = f(v, \sigma) \\
\sigma_S^{obs} = g(v, \sigma)
\end{cases} \quad \text{(9.81)}$$

\subsubsection{Iterative Approach ( KMV-Moody's Method) }

\textbf{Step 1:} Set initial asset volatility $\hat{\sigma}_0 = \sigma_S$ (equity volatility), then invert $f(v, \sigma)$ to obtain daily asset values:
$$v_j^{(0)} = f^{-1}(S_j^{obs}, \hat{\sigma}_0), \quad j = 1, \ldots, m. \quad \text{(9.82)}$$

\textbf{Step 2:} Update asset volatility by computing the standard deviation of log-returns $\{\ln(v_{j+1}/v_j)\}$ to get $\hat{\sigma}_1$, then repeat these two steps untill convergence :
$$v_j^{(1)} = f^{-1}(S_j^{obs}, \hat{\sigma}_1), \quad j = 1, \ldots, m. \quad \text{(9.83)}$$

Because of leverage effect, asset volatility converges from above.

\subsubsection{Maximum Likelihood }

Since the asset value process is a GBM, one-period log-returns under the physical measure are normally distributed: $\ln(v_{j+1}/v_j) \sim \mathcal{N}(\mu_d, \sigma_d^2)$. The joint likelihood of $m$ observations is:
$$L = \prod_{j=1}^m \frac{1}{\sigma_d\sqrt{2\pi}} \exp\left(-\frac{(\ln\frac{v_{j+1}}{v_j} - \mu_d)^2}{2\sigma_d^2}\right).$$

Maximizing the log-likelihood:
$$\arg\max_{\mu_d,\sigma_d} \ln L = \arg\max_{\mu_d,\sigma_d} \left(-\frac{1}{2}(m\ln 2\pi + m\ln\sigma_d^2) + \frac{1}{\sigma_d^2}\sum_{j=1}^m (\ln\frac{v_{j+1}}{v_j} - \mu_d)^2\right).$$

When asset value is unobservable, Duan (1994) shows that MLE can be performed on a contingent claim $f(v)$ (like market capitalization). Let $\hat{v}_j(\sigma)$ solve $S^{obs} = f(\hat{v}_j(\sigma), \sigma)$. The log-likelihood becomes:
$$\ln L = -\frac{1}{2}\left(m\ln 2\pi + m\ln\sigma_d^2 + \frac{1}{\sigma_d^2}\sum_{j=1}^m (\ln\frac{v_{j+1}}{v_j} - \mu_d)^2\right) - \sum_{j=1}^m \ln\left|\frac{\partial f}{\partial v}(\hat{v}_j(\sigma), \sigma)\right|.$$

The partial derivative $\partial f/\partial v$ is the Black-Scholes call delta. The MLE provides estimates for $\mu_d, \sigma_d$, and therefore $v$. When $v$ is driven by a GBM, the iterative method and MLE are asymptotically equivalent.

\subsection{Reminder: Binary Options}

\textbf{Cash-or-nothing:} Call $= M\Phi(d_2)e^{-rT}$, Put $= M\Phi(-d_2)e^{-rT}$; Payoffs: $\max(M \cdot \mathbf{1}_{v_T>K}, 0)$ and $\max(M \cdot \mathbf{1}_{v_T \leq K}, 0)$.

\textbf{Asset-or-nothing:} Call $= v\Phi(d_1)$, Put $= v\Phi(-d_1)$; Payoffs: $\max(v_T \cdot \mathbf{1}_{v_T>K}, 0)$ and $\max(v_T \cdot \mathbf{1}_{v_T \leq K}, 0)$.

$$d_1 = \frac{\ln(v/K) + (r + \sigma^2/2)T}{\sigma\sqrt{T}}, \quad d_2 = d_1 - \sigma\sqrt{T}$$


\section{Reduced-form Models}
View default as an exogenous event that can be statistically estimated or calibrated from market data.
therefore loses in terms of economic interpretation, but it relies on a weaker set of assumptions 
and makes full use of timely market information.

\subsection{ The Intensity-based Default Modeling }
\subsubsection{Preliminaries }
The cumulative survival function with constant hazard rate $h$:
$$\text{Proba}(\tau > t) = \exp(-ht)$$

More generally with time-dependent hazard rate:
$$\text{Proba}(\tau > t) = \exp\left(-\int_0^t h(s)ds\right)$$

For small time-step $\Delta t$, the conditional default probability:
$$\text{Proba}(\tau \leq t + \Delta t | \tau > t) = \frac{\text{Proba}(\tau \leq t + \Delta t, \tau > t)}{\text{Proba}(\tau > t)} = \frac{\exp\left(-\int_0^t h(s)ds\right) - \exp\left(-\int_0^{t+\Delta t} h(s)ds\right)}{\exp\left(-\int_0^t h(s)ds\right)}$$

Simplifying:
$$\text{Proba}(\tau \leq t + \Delta t | \tau > t) = 1 - \exp\left(-\int_t^{t+\Delta t} h(s)ds\right)$$

In the limit as $\Delta t \to 0$:
$$\lim_{\Delta t \to 0} \text{Proba}(\tau \leq t + \Delta t | \tau > t) = h(t)\Delta t.$$

In other words, $h(t)\Delta t$ is the conditional probability of default right after time $t$ given survival up to time $t$.

For pricing purposes, we reason under the EMM $Q$. The goal is to define a process $\{\lambda_t\}_{t \geq 0}$ such that:
$$Q(\tau > t) = \mathbb{E}\left[\exp\left(-\int_0^t \lambda_s ds\right)\right].$$

The process $\lambda_t$, which in all generality can be stochastic, is referred to as the \textbf{default intensity}. It bears the same interpretation as the hazard rate of an instantaneous default probability, but under the EMM.

The value of the defaultable zero-coupon bond (no recovery):
$$D(t,T) = \mathbb{E}_t\left[\exp\left(-\int_t^T (r_s + \lambda_s)ds\right)\right].$$

When the bond offers proportional recovery $\delta$ upon default (recovery of market value, RMV):
$$D(t,T) = \mathbb{E}_t\left[\exp\left(-\int_t^T (r_s + (1-\delta)\lambda_s)ds\right)\right].$$

\subsubsection{Parametric Specifications }
negative correlation between credit spreads and the level of risk-free rates is a strong stylized fact. That 
being said, both the instantaneous risk-free rate and the default intensity should be non-negative 
processes. The modelling challenge therefore consists in handling a negative correlation between 
two processes that must remain positive.

\textbf{Duffee (1999) Affine Model:} Assumes an affine relation between $\lambda_t$ and $r_t$ (with $dW_t dZ_t = 0$):
$$dr_t = \kappa_r(\theta_r - r_t)dt + \sigma_r\sqrt{r_t}dW_t, \quad ds_t = \kappa_s(\theta_s - s_t)dt + \sigma_s\sqrt{s_t}dZ_t,$$

with default intensity:
$$\lambda_t = \alpha + \beta r_t + s_t,$$

where $\alpha, \beta$ are constants. $\beta < 0$ captures negative correlation, though this doesn't guarantee $\lambda_t \geq 0$ always.

\subsubsection{Hybrid Models }
Instead of using a purely exogenous process to model the default intensity, one can relate the 
process $\lambda_t$ to factors that have been shown to be empirical determinants(factors) of credit spreads. some factors and their impact on credit spreads are listed below:

\noindent\textbf{Market factors:} $\Delta$ Level of yield curve $(-)$, $\Delta$ Slope of yield curve $(+)$, $\Delta$ Industrial production $(-)$, $\Delta$ VIX $(+)$, $\Delta$ Slope of volatility smirk $(+)$, $\Delta$ HML Fama-French factor $(-)$, $\Delta$ SMB Fama-French factor $(-)$, Return on S\&P 500 index $(-)$.

\noindent\textbf{Default:} $\Delta$ Leverage $(+)$, $\Delta$ Default probability $(+)$.

\noindent\textbf{Liquidity:} $\Delta$ PCA components $(+)$.

\textbf{Madan and Unal (2000) Model:} Example of a hybrid model linking intensity to firm factors:
$$\lambda_t = a + b\ln V_t + cr_t,$$
where $V_t$ is the firm's cash assets level. The model uses a Vasicek process for the risk-free rate and a correlated GBM for cash assets.


\subsection{Rating-based Pricing  }
\subsubsection{Rating Transition Matrices}
These matrices contain the historical frequencies of transitioning from one rating 
category to another. As expected, the mass of historical frequencies is concentrated on the diagonal (no change in the 
rating). remaining mass of frequencies is located for the most part on the 
adjacent categories, indicating a limited dispersion and a tendency for mean reversion. Finally, a 
significant number of cells is equal to zero. Large shifts in rating (including default) are rare events.

Let $p_{ij}$ denote the probability of switching from rating $i$ to rating $j$ within a year. Let $P$ denote the matrix containing all the $p_{ij}$. Under the assumptions of \textbf{time-invariance} (the transition probability from date $t$ to date $T$ depends on $T-t$ only) and \textbf{no memory} (the transition probability from rating $i$ to rating $j$ depends on $i$ and $j$ only without momentum effects), the rating transition matrix can be viewed as a \textbf{time-homogenous Markov chain}.

The $n$-year transition matrix is simply:
$$P^{(n)} = P^n.$$

Two-period transition probabilities: $p_{ij}^{(2)} = \sum_{k=1}^K p_{ik} \times p_{kj}$ with two-year transition matrix $P^{(2)} = P \times P$.

\textbf{Limitations of using historical rating transition matrices:}
\begin{itemize}
\item Using empirical transitions to compute default probabilities generates spreads that are too low for low maturities and too high for high maturities, indicating that transitions are not time-homogenous.
\item More fundamentally, empirical transitions are physical probabilities while bond pricing requires EMM probabilities.
\end{itemize}

\subsubsection{Bond Pricing }

\textbf{Objective:} Adjust historical transition matrix $P$ into time-nonhomogeneous matrix $Q$ containing EMM probabilities $q_{ij}$.

\textbf{RTV Recovery Assumption:} Defaultable zero-coupon bond in rating category $i$:
$$D_i(t,T) = P(t,T)(\delta + (1-\delta)Q_i(\tau > T | \tau > t)).$$

\textbf{Back out implied EMM default probability:} From observed bond prices,
$$q_i^{imp}(\tau < T) = \frac{P(0,T) - D_i(0,T)}{(1-\delta)P(0,T)}.$$

\textbf{Rating-specific risk premia:} Let $U$ be the diagonal matrix of risk premia $\mu_i$, k is the default column, then:
$$q_{ij} = \begin{cases} p_{ij}\mu_i & \text{for } j \neq K \\ 1 - \mu_i(1-p_{iK}) & \text{for } j = K \end{cases}$$

For first maturity: $\mu_i(1) = \frac{D_i(0,1) - \delta P(0,1)}{(1-\delta)P(0,1)(1-p_{iK})}$.

\subsection{Modeling the Recovery Rate}

Recovery rates show substantial variations across debt types, industries, and business conditions, with negative relations to default rates and high yield bond supply, and positive relations to GDP growth and S\&P 500 returns. This exhibits \textbf{procyclicality of recovery rates}.

\subsubsection{The Joint Identification Problem}

The \textbf{default probability and the recovery rate can be separately estimated if we observe the prices of two defaultable instruments issued by the same company}. From the \textit{pari passu} clause, default affects both instruments simultaneously, but their payoffs upon default differ.

Unal, Madan and Güntay (2003) extract recovery rates from a sample of firms with senior and junior debts. Assuming an RTV recovery rate that is independent of the default timing:
$$\frac{D_s(0,T) - D_j(0,T)}{P(0,T) - D_j(0,T)} = \frac{\mathbb{E}(\delta_s) - \mathbb{E}(\delta_j)}{1 - \mathbb{E}(\delta_j)}.$$

Das and Hanouna (2009) develop a method to extract default probabilities and recovery rates from the joint observation of equity and CDS prices, using a binomial tree where equity prices drive default intensity (negatively correlated), and recovery rates follow a logistic function of intensity. Under absolute priority, shareholders get nothing at default while CDS recovery is logistic in intensity.

\textbf{Recovery term structure and dynamics:} The term structure is downward sloping. Firms experiencing sudden default obtain greater resale value for assets, while firms whose financial health deteriorates gradually see assets lose resale value. Recovery rates are increasing with credit rating quality (CCC near default shows substantially lower recoveries). Significant variations in recovery exist among industries.

\subsubsection{The Beta Distribution}

Once recovery rates are empirically extracted, their distribution can be smoothed with a probability law. The beta distribution is ideal because it is defined on $[0,1]$ and is parsimonious (two-parameter), producing diverse distribution shapes.

For $x \in [0,1]$, the beta density is:
$$f(x;a,b) = \frac{x^{a-1}(1-x)^{b-1}}{B(a,b)},$$
where $a, b > 0$ are parameters and $B(a,b) = \frac{\Gamma(a)\Gamma(b)}{\Gamma(a+b)}$ is the scaling constant. Depending on $a$ and $b$, the beta distribution can be increasing, decreasing, bell-shaped, U-shaped, right-skewed, or left-skewed.

The mean is $a/(a+b)$ and variance is $ab/[(a+b)^2(a+b+1)]$. To match empirical mean $\mu$ and standard deviation $\sigma$ of recovery rates, the first two moments are matched:
$$a = \frac{\mu^2 - \mu^3 - \mu\sigma^2}{\sigma^2}, \quad b = \frac{\mu - 2\mu^2 + \mu^3 - \sigma^2 + \mu\sigma^2}{\sigma^2}.$$

When $a$ and $b$ are close to small integers, a simple simulation: generate $a+b-1$ uniform draws from $[0,1]$ and select the $a$-th lowest. For non small-integer parameters, use gamma distribution draws.

Renault and Scaillet (2004) improve fitting with a beta kernel estimator. Given a sample of $n$ recovery rates $\delta_1, \ldots, \delta_n$, the estimator is:
$$\hat{f}(x) = \frac{1}{n}\sum_{j=1}^n f\left(\delta_j; 1 + \frac{x}{w}, 1 + \frac{1-x}{w}\right),$$
where $w$ is the bandwidth smoothing parameter. Renault and Scaillet (2004) set $w = \sigma(n)^{-2/5}$ with $\sigma$ the empirical standard deviation.

While the beta distribution provides an acceptable approximation, the kernel beta estimator reveals that recovery rates can be bimodal for some instruments. However, this bimodality is observed in empirical data (physical probabilities). In contrast, recovery rates implied from market prices (bonds or CDS) typically exhibit a unimodal distribution.

\end{document}
